
\documentclass[12pt,a4paper]{article}

\usepackage[T2A]{fontenc}
\usepackage[utf8]{inputenc}
\usepackage[russian]{babel}
\usepackage{amsmath,amssymb,amsthm}
\usepackage{enumitem}
\usepackage{array}
\usepackage{booktabs}
\usepackage{mathtools}
\usepackage{geometry}
\usepackage{hyperref}

\geometry{left=2.2cm,right=2.2cm,top=2cm,bottom=2cm}

\newtheorem{definition}{Определение}
\newtheorem{theorem}{Теорема}
\newtheorem{remark}{Замечание}
\newtheorem{example}{Пример}

\title{Билет №26 \\ \small Особенности архитектуры локальных сетей (Ethernet, FDDI, WiFi). (ИТМО) \\ Особенности физической архитектуры локальных сетей: Ethernet (коаксиальные кабели, витые пары, оптоволоконные сети), Token Ring, FDDI. (СпБГУ)}
\author{Сериков Владислав Александрович}
\date{6 мая 2026}

\begin{document}

\maketitle
\tableofcontents
\newpage

\section{Архитектура локальных сетей: Ethernet, Token Ring, FDDI, Wi-Fi}

\subsection{Базовые понятия локальных сетей}

\textbf{Локальная вычислительная сеть} (LAN, Local Area Network) --- это сеть передачи данных, обычно охватывающая ограниченную территорию: аудиторию, лабораторию, здание, кампус, офис или производственный участок. Для локальных сетей характерны:
\begin{itemize}
    \item относительно малые задержки;
    \item высокая пропускная способность;
    \item административная принадлежность одной организации или ограниченному числу организаций;
    \item использование канального и физического уровней семейства стандартов IEEE 802.
\end{itemize}

Архитектуру локальной сети удобно рассматривать в нескольких аспектах:
\begin{enumerate}
    \item \textbf{Физическая топология} --- реальное соединение узлов кабелями, радиоканалами, концентраторами, коммутаторами, точками доступа.
    \item \textbf{Логическая топология} --- способ, которым узлы воспринимают среду передачи данных: общая шина, кольцо, звезда, дерево, ячеистая структура.
    \item \textbf{Метод доступа к среде} --- правило, определяющее, какой узел и когда имеет право передавать кадр.
    \item \textbf{Канальный формат кадров} --- структура MAC-кадров, адресация, контроль ошибок.
    \item \textbf{Физическая среда} --- коаксиальный кабель, витая пара, оптоволокно, радиоканал.
\end{enumerate}

В модели OSI локальные сети IEEE 802 в основном описывают два нижних уровня:
\begin{itemize}
    \item \textbf{физический уровень} --- электрические, оптические или радиосигналы, кодирование, разъёмы, частоты, скорости;
    \item \textbf{канальный уровень} --- формирование кадров, MAC-адресация, доступ к среде, обнаружение ошибок.
\end{itemize}

Канальный уровень в IEEE 802 часто делят на два подуровня:
\begin{itemize}
    \item \textbf{LLC} (Logical Link Control) --- логическое управление каналом, интерфейс к сетевому уровню;
    \item \textbf{MAC} (Medium Access Control) --- управление доступом к среде передачи.
\end{itemize}

\subsection{Классификация топологий локальных сетей}

\subsubsection{Физическая шина}

При шинной топологии все станции подключаются к общей передающей среде.

\[
\text{Узел}_1 \;-\; \text{Узел}_2 \;-\; \text{Узел}_3 \;-\; \text{Узел}_4
\]

Особенности:
\begin{itemize}
    \item все узлы разделяют одну среду;
    \item в каждый момент времени корректно может передавать только один узел;
    \item необходим механизм разрешения конфликтов;
    \item повреждение магистрального кабеля может нарушить работу всей сети.
\end{itemize}

Исторический пример: ранний Ethernet на коаксиальном кабеле: 10BASE5 и 10BASE2.

\subsubsection{Физическая звезда}

Все станции подключаются к центральному устройству: концентратору, коммутатору или точке доступа.

\[
\begin{array}{ccccc}
 & & \text{Коммутатор} & & \\
 & \diagup & | & \diagdown & \\
\text{Узел}_1 & & \text{Узел}_2 & & \text{Узел}_3
\end{array}
\]

Особенности:
\begin{itemize}
    \item отказ отдельного кабеля обычно отключает только одну станцию;
    \item центральное устройство является важной точкой отказа;
    \item удобно масштабировать и администрировать;
    \item современный Ethernet на витой паре физически строится как звезда.
\end{itemize}

\subsubsection{Логическое кольцо}

В логическом кольце право передачи последовательно переходит от одной станции к другой.

\[
\text{Узел}_1 \rightarrow \text{Узел}_2 \rightarrow \text{Узел}_3 \rightarrow \text{Узел}_4 \rightarrow \text{Узел}_1
\]

Особенности:
\begin{itemize}
    \item передача обычно детерминирована;
    \item можно гарантировать верхнюю границу времени ожидания доступа к среде;
    \item сложнее реализовать отказоустойчивость;
    \item примеры: Token Ring, FDDI.
\end{itemize}

\subsubsection{Дерево и иерархическая звезда}

Современные Ethernet-сети часто имеют иерархическую структуру:

\[
\text{Ядро} \rightarrow \text{Распределение} \rightarrow \text{Доступ}
\]

На уровне доступа подключаются конечные устройства, на уровне распределения агрегируются потоки, на уровне ядра обеспечивается высокоскоростная коммутация между крупными сегментами.

\subsection{Канальный кадр и MAC-адресация}

\textbf{MAC-адрес} --- аппаратный адрес сетевого интерфейса канального уровня. Обычно имеет длину 48 бит и записывается шестью октетами в шестнадцатеричном виде:

\[
\texttt{00:1A:2B:3C:4D:5E}.
\]

Кадр канального уровня обычно содержит:
\begin{itemize}
    \item адрес получателя;
    \item адрес отправителя;
    \item служебные поля;
    \item поле данных;
    \item контрольную сумму кадра.
\end{itemize}

Для обнаружения ошибок в Ethernet и многих других технологиях используется поле FCS, основанное на циклическом избыточном коде CRC.

\subsection{Ethernet}

\subsubsection{Общая характеристика Ethernet}

\textbf{Ethernet} --- семейство технологий локальных сетей, стандартизованных в IEEE 802.3. Первоначально Ethernet был технологией общей шины с конкурентным доступом к среде, но современный Ethernet почти всегда строится на коммутаторах и полнодуплексных соединениях.

Основные свойства Ethernet:
\begin{itemize}
    \item использование MAC-адресов длиной 48 бит;
    \item кадрный способ передачи данных;
    \item исторически --- метод доступа CSMA/CD;
    \item современные реализации --- коммутируемая полнодуплексная передача без коллизий;
    \item поддержка различных физических сред: коаксиальный кабель, витая пара, оптоволокно;
    \item широкий диапазон скоростей: от 10 Мбит/с до сотен Гбит/с и выше.
\end{itemize}

\subsubsection{Формат кадра Ethernet II}

На практике широко используется формат Ethernet II:

\[
\begin{array}{|c|c|c|c|c|}
\hline
\text{Преамбула} & \text{SFD} & \text{MAC получателя} & \text{MAC отправителя} & \text{Тип} \\
\hline
7 \text{ байт} & 1 \text{ байт} & 6 \text{ байт} & 6 \text{ байт} & 2 \text{ байта} \\
\hline
\end{array}
\]

\[
\begin{array}{|c|c|}
\hline
\text{Данные} & \text{FCS} \\
\hline
46\text{--}1500 \text{ байт} & 4 \text{ байта} \\
\hline
\end{array}
\]

Назначение полей:
\begin{itemize}
    \item \textbf{Преамбула} --- синхронизация приёмника с передатчиком.
    \item \textbf{SFD} (Start Frame Delimiter) --- признак начала кадра.
    \item \textbf{MAC получателя} --- адрес назначения.
    \item \textbf{MAC отправителя} --- адрес источника.
    \item \textbf{Тип} --- идентификатор протокола верхнего уровня, например IPv4 или ARP.
    \item \textbf{Данные} --- полезная нагрузка.
    \item \textbf{FCS} --- контрольная последовательность кадра.
\end{itemize}

Минимальный размер Ethernet-кадра от MAC-адреса получателя до FCS равен:

\[
6 + 6 + 2 + 46 + 4 = 64 \text{ байта}.
\]

Максимальный стандартный размер без преамбулы и SFD:

\[
6 + 6 + 2 + 1500 + 4 = 1518 \text{ байт}.
\]

Если используется VLAN-тег IEEE 802.1Q, добавляются 4 байта, и размер кадра становится 1522 байта.

\subsubsection{Физическая архитектура Ethernet на коаксиальном кабеле}

Ранний Ethernet использовал коаксиальный кабель и физическую топологию общей шины.

\paragraph{10BASE5}

10BASE5, или ``толстый Ethernet'', имел следующие характеристики:
\begin{itemize}
    \item скорость передачи: 10 Мбит/с;
    \item среда: толстый коаксиальный кабель;
    \item физическая топология: шина;
    \item максимальная длина сегмента: порядка 500 м;
    \item подключение станций через трансиверы и AUI-интерфейс;
    \item на концах кабеля требовались терминаторы.
\end{itemize}

Схема:

\[
\text{Терминатор} - \text{Узел}_1 - \text{Узел}_2 - \text{Узел}_3 - \text{Терминатор}
\]

\paragraph{10BASE2}

10BASE2, или ``тонкий Ethernet'', использовал более тонкий коаксиальный кабель.

Характеристики:
\begin{itemize}
    \item скорость: 10 Мбит/с;
    \item среда: тонкий коаксиальный кабель;
    \item физическая топология: шина;
    \item максимальная длина сегмента: порядка 185 м;
    \item подключение через BNC T-коннекторы;
    \item требовались терминаторы на концах сегмента.
\end{itemize}

Преимущества по сравнению с 10BASE5:
\begin{itemize}
    \item дешевле;
    \item проще монтаж;
    \item не требуются внешние трансиверы.
\end{itemize}

Недостатки:
\begin{itemize}
    \item разрыв кабеля или плохой контакт мог нарушить работу всего сегмента;
    \item сложная диагностика;
    \item все станции находились в одном домене коллизий.
\end{itemize}

\subsubsection{Ethernet на витой паре}

С переходом к витой паре физическая топология Ethernet стала звездой. Каждый узел подключается отдельным кабелем к концентратору или коммутатору.

\[
\begin{array}{ccccc}
 & & \text{Коммутатор} & & \\
 & \diagup & | & \diagdown & \\
\text{ПК}_1 & & \text{ПК}_2 & & \text{ПК}_3
\end{array}
\]

Основные варианты:
\begin{itemize}
    \item \textbf{10BASE-T} --- 10 Мбит/с;
    \item \textbf{100BASE-TX} --- 100 Мбит/с;
    \item \textbf{1000BASE-T} --- 1 Гбит/с;
    \item \textbf{10GBASE-T} --- 10 Гбит/с.
\end{itemize}

Особенности Ethernet на витой паре:
\begin{itemize}
    \item используется кабель UTP или STP;
    \item типичная максимальная длина горизонтального кабельного участка --- 100 м;
    \item используется разъём 8P8C, часто называемый RJ-45;
    \item в современных сетях соединения обычно полнодуплексные;
    \item каждый порт коммутатора образует отдельный коллизионный домен.
\end{itemize}

\paragraph{Концентратор и коммутатор}

\textbf{Концентратор} (hub) работает на физическом уровне. Он повторяет сигнал со входного порта на все остальные порты. Все устройства, подключённые к концентратору, находятся в одном домене коллизий.

\textbf{Коммутатор} (switch) работает преимущественно на канальном уровне. Он анализирует MAC-адреса и пересылает кадры только в нужные порты.

Пример таблицы коммутации:

\[
\begin{array}{|c|c|}
\hline
\text{MAC-адрес} & \text{Порт} \\
\hline
00:AA:00:00:00:01 & 1 \\
00:AA:00:00:00:02 & 2 \\
00:AA:00:00:00:03 & 3 \\
\hline
\end{array}
\]

\subsubsection{Ethernet на оптоволокне}

Оптоволоконный Ethernet используется для:
\begin{itemize}
    \item магистральных соединений;
    \item соединения зданий;
    \item высокоскоростных каналов между коммутаторами;
    \item сетей центров обработки данных;
    \item участков с сильными электромагнитными помехами.
\end{itemize}

Основные преимущества оптоволокна:
\begin{itemize}
    \item высокая пропускная способность;
    \item большие расстояния передачи;
    \item нечувствительность к электромагнитным помехам;
    \item электрическая развязка между устройствами;
    \item сложность несанкционированного подключения по сравнению с медным кабелем.
\end{itemize}

Типы волокон:
\begin{itemize}
    \item \textbf{многомодовое волокно} --- используется на сравнительно малых расстояниях, например внутри здания или ЦОД;
    \item \textbf{одномодовое волокно} --- используется на больших расстояниях.
\end{itemize}

Примеры стандартов:
\begin{itemize}
    \item \textbf{100BASE-FX} --- Fast Ethernet по оптоволокну;
    \item \textbf{1000BASE-SX} --- Gigabit Ethernet по многомодовому волокну;
    \item \textbf{1000BASE-LX} --- Gigabit Ethernet по одномодовому или многомодовому волокну;
    \item \textbf{10GBASE-SR} --- 10 Gigabit Ethernet на короткие расстояния по многомодовому волокну;
    \item \textbf{10GBASE-LR} --- 10 Gigabit Ethernet на большие расстояния по одномодовому волокну.
\end{itemize}

\subsubsection{Метод доступа CSMA/CD}

В классическом полудуплексном Ethernet использовался метод CSMA/CD.

\textbf{CSMA/CD} --- Carrier Sense Multiple Access with Collision Detection, множественный доступ с контролем несущей и обнаружением коллизий.

Идея:
\begin{itemize}
    \item перед передачей станция слушает среду;
    \item если среда свободна, станция начинает передачу;
    \item если две станции начали передачу почти одновременно, возникает коллизия;
    \item станции обнаруживают коллизию, прекращают передачу и повторяют попытку через случайный интервал.
\end{itemize}

\paragraph{Алгоритм CSMA/CD}

Пусть станция должна передать кадр.

\begin{enumerate}
    \item Станция проверяет наличие несущей в среде.
    \item Если среда занята, станция ждёт, пока среда освободится.
    \item Если среда свободна, станция начинает передачу кадра.
    \item Во время передачи станция продолжает контролировать среду.
    \item Если коллизия не обнаружена, кадр считается успешно переданным.
    \item Если обнаружена коллизия:
    \begin{enumerate}
        \item станция передаёт jam-сигнал, чтобы все участники заметили коллизию;
        \item прекращает передачу кадра;
        \item увеличивает счётчик попыток;
        \item выбирает случайную задержку по алгоритму двоичной экспоненциальной отсрочки;
        \item после задержки повторяет попытку.
    \end{enumerate}
    \item Если число попыток превысило допустимый предел, передача завершается ошибкой.
\end{enumerate}

\paragraph{Двоичная экспоненциальная отсрочка}

После $k$-й коллизии станция выбирает случайное число:

\[
r \in \{0,1,\ldots,2^m - 1\},
\]

где

\[
m = \min(k, 10).
\]

Задержка равна:

\[
T_{\text{backoff}} = r \cdot T_{\text{slot}}.
\]

Для классического Ethernet:

\[
T_{\text{slot}} = 512 \text{ битовых интервалов}.
\]

Для 10 Мбит/с:

\[
T_{\text{slot}} = \frac{512}{10^7} = 51{,}2 \ \mu s.
\]

\paragraph{Минимальный пример работы CSMA/CD}

Пусть станции $A$ и $B$ одновременно хотят передать кадры.

\[
A \longleftrightarrow \text{общая среда} \longleftrightarrow B
\]

\begin{enumerate}
    \item $A$ проверяет среду и видит, что она свободна.
    \item $B$ почти одновременно проверяет среду и тоже видит, что она свободна.
    \item $A$ и $B$ начинают передачу.
    \item Сигналы сталкиваются: возникает коллизия.
    \item Обе станции обнаруживают коллизию и передают jam-сигнал.
    \item $A$ выбирает, например, $r_A = 0$.
    \item $B$ выбирает, например, $r_B = 1$.
    \item $A$ повторяет передачу сразу после интервала ожидания.
    \item $B$ ждёт один slot time и затем повторяет попытку.
\end{enumerate}

В результате вероятность повторной коллизии уменьшается.

\subsubsection{Почему существует минимальный размер кадра Ethernet}

В полудуплексном Ethernet станция должна успеть обнаружить коллизию, пока она ещё передаёт кадр.

Пусть:
\begin{itemize}
    \item $\tau$ --- максимальное время распространения сигнала между двумя самыми удалёнными станциями;
    \item $2\tau$ --- время распространения сигнала туда и обратно;
    \item $T_f$ --- время передачи минимального кадра.
\end{itemize}

Для гарантированного обнаружения коллизии необходимо:

\[
T_f \geq 2\tau.
\]

Именно поэтому в классическом Ethernet минимальный кадр имеет длину 64 байта, то есть 512 бит.

\paragraph{Утверждение}

В полудуплексной сети с методом CSMA/CD передающая станция гарантированно обнаруживает любую коллизию тогда и только тогда, когда время передачи минимального кадра не меньше удвоенного максимального времени распространения сигнала между наиболее удалёнными станциями.

\[
T_f \geq 2\tau.
\]

\paragraph{Доказательство}

Рассмотрим две наиболее удалённые станции $A$ и $B$. Пусть максимальное время распространения сигнала между ними равно $\tau$.

\textbf{Необходимость.}

Пусть станция $A$ начала передачу в момент времени $t=0$. Станция $B$ находится на максимально возможном расстоянии от $A$.

Сигнал от $A$ достигнет $B$ только в момент времени $t=\tau$. Если $B$ начнёт передачу непосредственно перед приходом сигнала от $A$, то для $B$ среда ещё будет казаться свободной. Тогда произойдёт коллизия.

Информация о коллизии начнёт распространяться обратно к $A$ и достигнет её ещё через время $\tau$. Следовательно, $A$ узнает о коллизии примерно в момент:

\[
t = 2\tau.
\]

Если $A$ завершит передачу раньше, чем через $2\tau$, то она может не обнаружить коллизию. Поэтому необходимо:

\[
T_f \geq 2\tau.
\]

\textbf{Достаточность.}

Пусть теперь $T_f \geq 2\tau$. Рассмотрим любую возможную коллизию. Наихудший случай --- когда коллизия возникает у наиболее удалённой станции почти через $\tau$ после начала передачи первой станции. Сигнал о коллизии вернётся к первой станции не позднее чем ещё через $\tau$. Значит, первая станция обнаружит коллизию не позднее момента:

\[
t = 2\tau.
\]

Так как станция передаёт кадр не менее $T_f$, а $T_f \geq 2\tau$, она всё ещё будет находиться в состоянии передачи и сможет обнаружить коллизию.

Следовательно, условие достаточно.

Итак, условие

\[
T_f \geq 2\tau
\]

является необходимым и достаточным для гарантированного обнаружения коллизий в худшем случае.

\subsubsection{Современный коммутируемый Ethernet}

В современных Ethernet-сетях обычно используются коммутаторы и полнодуплексные каналы.

Особенности:
\begin{itemize}
    \item коллизии отсутствуют, так как передача и приём идут по разным направлениям или логически независимым каналам;
    \item CSMA/CD фактически не используется в полнодуплексном режиме;
    \item каждый порт коммутатора является отдельным сегментом;
    \item производительность сети значительно выше, чем у общей шины;
    \item возможны VLAN, агрегация каналов, QoS, Spanning Tree Protocol и другие механизмы.
\end{itemize}

\subsection{Token Ring}

\subsubsection{Общая характеристика Token Ring}

\textbf{Token Ring} --- технология локальных сетей, основанная на передаче специального маркера, или токена. Право передавать данные получает только та станция, которая владеет токеном.

Token Ring был стандартизован как IEEE 802.5 и активно использовался в корпоративных сетях до широкого распространения коммутируемого Ethernet.

Основные свойства:
\begin{itemize}
    \item логическая топология --- кольцо;
    \item физическая топология часто --- звезда с использованием MSAU;
    \item метод доступа --- передача маркера;
    \item детерминированный доступ к среде;
    \item типичные скорости классических реализаций: 4 и 16 Мбит/с.
\end{itemize}

\subsubsection{Физическая и логическая топология Token Ring}

Хотя логически Token Ring является кольцом, физически станции часто подключались к MSAU.

\[
\begin{array}{ccccc}
 & & \text{MSAU} & & \\
 & \diagup & | & \diagdown & \\
\text{Станция}_1 & & \text{Станция}_2 & & \text{Станция}_3
\end{array}
\]

Внутри MSAU соединения организуются так, что кадры проходят по логическому кольцу:

\[
\text{Станция}_1 \rightarrow \text{Станция}_2 \rightarrow \text{Станция}_3 \rightarrow \text{Станция}_1.
\]

MSAU может автоматически обходить отключённую или неисправную станцию, сохраняя целостность кольца.

\subsubsection{Метод доступа с передачей маркера}

\textbf{Маркер} --- короткий специальный кадр, циркулирующий по кольцу и предоставляющий право передачи.

\paragraph{Алгоритм работы Token Ring}

\begin{enumerate}
    \item По кольцу циркулирует свободный маркер.
    \item Станция, не имеющая данных для передачи, просто ретранслирует маркер дальше.
    \item Станция, имеющая данные, дожидается получения свободного маркера.
    \item Получив маркер, станция захватывает его и начинает передачу кадра данных.
    \item Кадр проходит по кольцу через промежуточные станции.
    \item Каждая промежуточная станция ретранслирует кадр дальше.
    \item Станция-получатель распознаёт свой адрес, копирует кадр и выставляет признаки получения.
    \item Кадр возвращается к отправителю.
    \item Отправитель удаляет свой кадр из кольца.
    \item Отправитель выпускает новый свободный маркер.
\end{enumerate}

\paragraph{Минимальный пример}

Пусть есть кольцо из трёх станций:

\[
A \rightarrow B \rightarrow C \rightarrow A.
\]

Станция $A$ хочет передать кадр станции $C$.

\begin{enumerate}
    \item Свободный маркер приходит к $A$.
    \item $A$ захватывает маркер.
    \item $A$ отправляет кадр с адресом получателя $C$.
    \item $B$ получает кадр, видит, что адрес не её, и передаёт дальше.
    \item $C$ получает кадр, копирует данные и отмечает факт приёма.
    \item Кадр возвращается к $A$.
    \item $A$ удаляет кадр из кольца.
    \item $A$ выпускает свободный маркер.
\end{enumerate}

\subsubsection{Преимущества и недостатки Token Ring}

Преимущества:
\begin{itemize}
    \item отсутствие коллизий;
    \item предсказуемое время доступа;
    \item хорошая работа при высокой нагрузке;
    \item возможность приоритетного доступа.
\end{itemize}

Недостатки:
\begin{itemize}
    \item сложность оборудования;
    \item высокая стоимость по сравнению с Ethernet;
    \item чувствительность к отказам кольцевой логики;
    \item меньшая распространённость;
    \item вытеснение коммутируемым Ethernet.
\end{itemize}

\subsection{FDDI}

\subsubsection{Общая характеристика FDDI}

\textbf{FDDI} (Fiber Distributed Data Interface) --- технология локальных и кампусных сетей, использующая оптоволоконную среду и маркерный доступ. Она была рассчитана на скорость 100 Мбит/с и часто применялась как магистральная технология до широкого распространения Fast Ethernet и Gigabit Ethernet.

Основные свойства:
\begin{itemize}
    \item скорость передачи: 100 Мбит/с;
    \item среда: преимущественно оптоволокно;
    \item логическая топология: двойное кольцо;
    \item метод доступа: передача маркера;
    \item высокая отказоустойчивость за счёт второго кольца;
    \item возможность охвата кампусных сетей.
\end{itemize}

\subsubsection{Двойное кольцо FDDI}

Классическая FDDI-сеть использует два кольца:
\begin{itemize}
    \item \textbf{первичное кольцо} --- основная передача данных;
    \item \textbf{вторичное кольцо} --- резервное кольцо, обычно передающее в противоположном направлении.
\end{itemize}

Схематически:

\[
\begin{array}{cccc}
A & \rightarrow & B & \rightarrow \\
\uparrow & & & \downarrow \\
D & \leftarrow & C & \leftarrow
\end{array}
\]

Вторичное кольцо используется для восстановления связности при отказе.

\subsubsection{Типы подключений FDDI}

В FDDI различают несколько типов станций:
\begin{itemize}
    \item \textbf{DAS} (Dual Attachment Station) --- станция с подключением к двум кольцам;
    \item \textbf{SAS} (Single Attachment Station) --- станция с подключением к одному кольцу через концентратор;
    \item \textbf{DAC} (Dual Attachment Concentrator) --- концентратор с двойным подключением;
    \item \textbf{SAC} (Single Attachment Concentrator) --- концентратор с одинарным подключением.
\end{itemize}

DAS подключается непосредственно к обоим кольцам:

\[
\text{Кольцо 1} \longrightarrow \boxed{\text{DAS}} \longrightarrow \text{Кольцо 1}
\]

\[
\text{Кольцо 2} \longleftarrow \boxed{\text{DAS}} \longleftarrow \text{Кольцо 2}
\]

SAS обычно подключается через концентратор, что упрощает администрирование и повышает устойчивость сети к отказам конечных станций.

\subsubsection{Восстановление FDDI при разрыве кольца}

Одно из важных свойств FDDI --- возможность автоматического восстановления при повреждении линии или отказе узла. При разрыве кольца соседние станции выполняют \textbf{wrap} --- замыкание первичного и вторичного колец в единую рабочую петлю.

\paragraph{Алгоритм восстановления FDDI при одиночном разрыве}

\begin{enumerate}
    \item Станции постоянно контролируют наличие сигнала и корректность прохождения кадров.
    \item При обнаружении разрыва соседние станции фиксируют потерю связи.
    \item Каждая из соседних станций переводит порт в режим обхода.
    \item Первичное кольцо соединяется со вторичным.
    \item Формируется новое логическое кольцо, обходя повреждённый участок.
    \item Передача данных продолжается по восстановленной топологии.
\end{enumerate}

\paragraph{Минимальный пример}

Изначально имеется кольцо:

\[
A \rightarrow B \rightarrow C \rightarrow D \rightarrow A.
\]

Пусть произошёл разрыв между $B$ и $C$.

Первичное кольцо нарушено:

\[
A \rightarrow B \;\; \times \;\; C \rightarrow D \rightarrow A.
\]

Станции $B$ и $C$ выполняют замыкание на вторичное кольцо. В результате образуется маршрут:

\[
A \rightarrow B \Rightarrow A \Leftarrow D \Leftarrow C,
\]

то есть повреждённый участок обходится по резервному направлению. Реальная схема зависит от конкретной реализации, но принцип состоит в соединении двух колец в одну непрерывную петлю.

\subsubsection{Метод доступа FDDI}

FDDI использует маркерный доступ, но с более сложным управлением временем, чем классический Token Ring.

Основные понятия:
\begin{itemize}
    \item \textbf{token} --- маркер, дающий право передачи;
    \item \textbf{TTRT} (Target Token Rotation Time) --- целевое время оборота маркера;
    \item \textbf{TRT} (Token Rotation Time) --- фактическое время оборота маркера;
    \item \textbf{synchronous traffic} --- синхронный трафик, обслуживаемый с гарантированной долей времени;
    \item \textbf{asynchronous traffic} --- асинхронный трафик, передаваемый при наличии оставшегося времени.
\end{itemize}

\paragraph{Упрощённый алгоритм доступа FDDI}

\begin{enumerate}
    \item Маркер циркулирует по кольцу.
    \item Станция измеряет время с момента предыдущего получения маркера.
    \item Если станция получает маркер, она сначала может передать синхронный трафик в пределах выделенной квоты.
    \item Затем, если фактический оборот маркера меньше допустимого целевого времени, станция может передавать асинхронный трафик.
    \item После исчерпания разрешённого времени передачи станция выпускает маркер дальше.
\end{enumerate}

Преимущество такого подхода состоит в сочетании:
\begin{itemize}
    \item предсказуемости;
    \item поддержки трафика с ограничениями по задержке;
    \item эффективного использования свободной полосы пропускания.
\end{itemize}

\subsubsection{Преимущества и недостатки FDDI}

Преимущества:
\begin{itemize}
    \item высокая для своего времени скорость 100 Мбит/с;
    \item оптоволоконная среда;
    \item большие расстояния;
    \item отказоустойчивость за счёт двойного кольца;
    \item предсказуемый доступ к среде.
\end{itemize}

Недостатки:
\begin{itemize}
    \item высокая стоимость оборудования;
    \item сложность монтажа и эксплуатации;
    \item вытеснение Ethernet-технологиями;
    \item избыточность для большинства современных локальных сетей.
\end{itemize}

\subsection{Wi-Fi}

\subsubsection{Общая характеристика Wi-Fi}

\textbf{Wi-Fi} --- семейство беспроводных локальных сетей на основе стандартов IEEE 802.11. В отличие от проводного Ethernet, Wi-Fi использует радиоканал, который является общей и менее предсказуемой средой передачи.

Основные особенности Wi-Fi:
\begin{itemize}
    \item физическая среда --- радиоканал;
    \item стандарты семейства IEEE 802.11;
    \item работа в нелицензируемых диапазонах, например 2{,}4 ГГц, 5 ГГц и 6 ГГц;
    \item поддержка мобильности клиентов;
    \item наличие точек доступа;
    \item использование CSMA/CA вместо CSMA/CD;
    \item полудуплексный характер радиоканала на одном канале.
\end{itemize}

\subsubsection{Архитектурные элементы Wi-Fi}

Основные элементы:
\begin{itemize}
    \item \textbf{STA} (station) --- станция, то есть устройство с Wi-Fi-интерфейсом;
    \item \textbf{AP} (access point) --- точка доступа;
    \item \textbf{BSS} (Basic Service Set) --- базовый набор обслуживания;
    \item \textbf{ESS} (Extended Service Set) --- расширенный набор обслуживания;
    \item \textbf{SSID} --- имя беспроводной сети;
    \item \textbf{DS} (Distribution System) --- распределительная система, обычно проводная Ethernet-сеть, соединяющая точки доступа.
\end{itemize}

\paragraph{Infrastructure BSS}

Наиболее распространённый режим --- инфраструктурный. Клиенты взаимодействуют через точку доступа.

\[
\begin{array}{ccccc}
\text{STA}_1 & \sim & \text{AP} & \sim & \text{STA}_2 \\
 & & | & & \\
 & & \text{Ethernet DS} & &
\end{array}
\]

Даже если две станции находятся рядом, в инфраструктурном режиме кадры обычно проходят через AP.

\paragraph{Ad Hoc / IBSS}

В режиме IBSS станции взаимодействуют напрямую без точки доступа:

\[
\text{STA}_1 \sim \text{STA}_2 \sim \text{STA}_3.
\]

Такой режим используется значительно реже.

\paragraph{ESS}

ESS объединяет несколько точек доступа с общим SSID:

\[
\text{AP}_1 \longleftrightarrow \text{DS} \longleftrightarrow \text{AP}_2.
\]

Клиент может перемещаться между зонами покрытия AP, выполняя роуминг.

\subsubsection{Почему в Wi-Fi используется CSMA/CA, а не CSMA/CD}

В Ethernet станция может передавать и одновременно контролировать среду, обнаруживая коллизии. В Wi-Fi это затруднено по причинам:
\begin{itemize}
    \item собственный передаваемый радиосигнал намного сильнее принимаемого;
    \item не все станции слышат друг друга;
    \item возможна проблема скрытого узла;
    \item радиоканал подвержен помехам и затуханию.
\end{itemize}

Поэтому Wi-Fi использует не обнаружение коллизий, а их предотвращение.

\textbf{CSMA/CA} --- Carrier Sense Multiple Access with Collision Avoidance, множественный доступ с контролем несущей и предотвращением коллизий.

\subsubsection{Алгоритм CSMA/CA в Wi-Fi}

Упрощённый алгоритм передачи кадра:

\begin{enumerate}
    \item Станция проверяет, свободен ли радиоканал.
    \item Если канал занят, станция откладывает передачу.
    \item Если канал свободен в течение межкадрового интервала DIFS, станция выбирает случайный backoff-счётчик.
    \item Пока канал свободен, счётчик уменьшается.
    \item Если канал становится занятым, счётчик замораживается.
    \item Когда канал снова свободен в течение DIFS, отсчёт продолжается.
    \item Когда счётчик достигает нуля, станция передаёт кадр.
    \item Получатель при успешном приёме отправляет ACK после короткого интервала SIFS.
    \item Если ACK не получен, станция считает передачу неудачной и повторяет попытку с увеличенным окном конкуренции.
\end{enumerate}

\paragraph{Минимальный пример CSMA/CA}

Пусть станции $A$ и $B$ хотят передать кадры через одну точку доступа.

\[
A \sim \text{AP} \sim B
\]

\begin{enumerate}
    \item Обе станции видят, что канал свободен.
    \item $A$ выбирает случайный backoff $3$ слота.
    \item $B$ выбирает случайный backoff $7$ слотов.
    \item Канал остаётся свободным.
    \item Счётчики уменьшаются:
    \[
    A: 3 \rightarrow 2 \rightarrow 1 \rightarrow 0,
    \]
    \[
    B: 7 \rightarrow 6 \rightarrow 5 \rightarrow 4.
    \]
    \item $A$ начинает передачу.
    \item $B$ слышит занятую среду и замораживает счётчик на значении $4$.
    \item AP принимает кадр $A$ и отправляет ACK.
    \item После освобождения канала $B$ продолжает отсчёт:
    \[
    4 \rightarrow 3 \rightarrow 2 \rightarrow 1 \rightarrow 0.
    \]
    \item $B$ передаёт свой кадр.
\end{enumerate}

\subsubsection{Проблема скрытого узла}

Пусть станции $A$ и $C$ не слышат друг друга, но обе слышат точку доступа $B$.

\[
A \sim B \sim C,
\]

но

\[
A \not\sim C.
\]

Если $A$ передаёт кадр точке доступа, станция $C$ может не обнаружить передачу $A$ и тоже начать передачу. В точке доступа возникнет коллизия.

Для уменьшения этой проблемы может использоваться механизм RTS/CTS.

\subsubsection{RTS/CTS}

RTS/CTS --- дополнительный механизм резервирования среды:
\begin{itemize}
    \item \textbf{RTS} (Request To Send) --- запрос на передачу;
    \item \textbf{CTS} (Clear To Send) --- разрешение на передачу.
\end{itemize}

\paragraph{Алгоритм RTS/CTS}

\begin{enumerate}
    \item Станция $A$ хочет передать большой кадр точке доступа $B$.
    \item $A$ передаёт короткий кадр RTS.
    \item $B$ отвечает кадром CTS.
    \item Все станции, услышавшие RTS или CTS, обновляют свой NAV --- Network Allocation Vector.
    \item Пока NAV положителен, станции считают среду занятой и не передают.
    \item $A$ передаёт основной кадр данных.
    \item $B$ подтверждает приём кадром ACK.
\end{enumerate}

\paragraph{Минимальный пример RTS/CTS}

\[
A \sim B \sim C, \qquad A \not\sim C.
\]

\begin{enumerate}
    \item $A$ отправляет RTS к $B$.
    \item $B$ отправляет CTS.
    \item $C$ слышит CTS от $B$ и понимает, что среда зарезервирована.
    \item $C$ временно не передаёт.
    \item $A$ передаёт данные к $B$.
    \item $B$ отправляет ACK.
\end{enumerate}

Так снижается вероятность коллизии у точки доступа.

\subsubsection{Физический уровень Wi-Fi}

Физический уровень Wi-Fi зависит от конкретной версии IEEE 802.11. Развитие стандартов включало:
\begin{itemize}
    \item 802.11b --- 2{,}4 ГГц, до 11 Мбит/с;
    \item 802.11a --- 5 ГГц, до 54 Мбит/с;
    \item 802.11g --- 2{,}4 ГГц, до 54 Мбит/с;
    \item 802.11n --- MIMO, работа в 2{,}4 и 5 ГГц;
    \item 802.11ac --- высокая скорость в диапазоне 5 ГГц;
    \item 802.11ax --- Wi-Fi 6/6E, OFDMA, улучшенная работа в плотных сетях;
    \item 802.11be --- Wi-Fi 7, дальнейшее увеличение скорости и эффективности.
\end{itemize}

Ключевые физические технологии:
\begin{itemize}
    \item \textbf{OFDM} --- ортогональное частотное мультиплексирование;
    \item \textbf{MIMO} --- использование нескольких передающих и принимающих антенн;
    \item \textbf{OFDMA} --- разделение канала между несколькими пользователями по поднесущим;
    \item \textbf{beamforming} --- формирование направленного радиолуча;
    \item \textbf{каналы различной ширины} --- 20, 40, 80, 160 МГц и др.
\end{itemize}

\subsubsection{Особенности беспроводной среды}

Wi-Fi существенно отличается от проводного Ethernet:
\begin{itemize}
    \item среда общая и разделяемая всеми станциями на одном канале;
    \item пропускная способность зависит от расстояния, помех, стен, отражений;
    \item скорость соединения может динамически изменяться;
    \item безопасность требует криптографической защиты;
    \item фактическая полезная скорость ниже номинальной физической скорости;
    \item возможны коллизии, скрытые узлы, интерференция с соседними сетями.
\end{itemize}

\subsection{Сравнение Ethernet, Token Ring, FDDI и Wi-Fi}

\[
\begin{array}{|p{3cm}|p{3cm}|p{3cm}|p{3cm}|}
\hline
\textbf{Технология} & \textbf{Топология} & \textbf{Метод доступа} & \textbf{Среда} \\
\hline
Ethernet ранний & Шина & CSMA/CD & Коаксиальный кабель \\
\hline
Ethernet современный & Звезда / дерево & Коммутация, full-duplex & Витая пара, оптоволокно \\
\hline
Token Ring & Логическое кольцо, физическая звезда & Передача маркера & Витая пара, STP/UTP \\
\hline
FDDI & Двойное кольцо & Передача маркера & Оптоволокно \\
\hline
Wi-Fi & Инфраструктурная BSS/ESS, ad hoc & CSMA/CA & Радиоканал \\
\hline
\end{array}
\]

\subsection{Сравнение методов доступа к среде}

\subsubsection{CSMA/CD}

Характерно для классического Ethernet.

\begin{itemize}
    \item Узлы соревнуются за среду.
    \item Коллизии возможны.
    \item Коллизии обнаруживаются во время передачи.
    \item После коллизии используется случайная экспоненциальная задержка.
    \item Хорошо работает при малой и средней нагрузке.
    \item В современном полнодуплексном Ethernet практически не применяется.
\end{itemize}

\subsubsection{Передача маркера}

Характерна для Token Ring и FDDI.

\begin{itemize}
    \item Передавать может только владелец маркера.
    \item Коллизии отсутствуют.
    \item Доступ детерминирован.
    \item При высокой нагрузке работает предсказуемо.
    \item Требуется управление маркером и восстановление при его потере.
\end{itemize}

\subsubsection{CSMA/CA}

Характерно для Wi-Fi.

\begin{itemize}
    \item Станции слушают канал перед передачей.
    \item Коллизии не обнаруживаются непосредственно, а предотвращаются вероятностно.
    \item Используются случайные backoff-интервалы.
    \item Успешная передача подтверждается ACK.
    \item Может применяться RTS/CTS.
\end{itemize}

\subsection{Физические среды локальных сетей}

\subsubsection{Коаксиальный кабель}

Коаксиальный кабель состоит из центрального проводника, диэлектрика, экрана и внешней оболочки.

\[
\text{проводник} \;|\; \text{диэлектрик} \;|\; \text{экран} \;|\; \text{оболочка}
\]

Преимущества:
\begin{itemize}
    \item хорошая помехозащищённость;
    \item возможность построения общей шины;
    \item исторически был удобен для ранних Ethernet-сетей.
\end{itemize}

Недостатки:
\begin{itemize}
    \item неудобство монтажа;
    \item сложная диагностика обрывов;
    \item низкая гибкость масштабирования;
    \item вытеснен витой парой и оптоволокном.
\end{itemize}

\subsubsection{Витая пара}

Витая пара состоит из пар медных проводников, скрученных для уменьшения электромагнитных наводок.

Типы:
\begin{itemize}
    \item \textbf{UTP} --- неэкранированная витая пара;
    \item \textbf{STP/FTP} --- экранированные варианты.
\end{itemize}

Преимущества:
\begin{itemize}
    \item низкая стоимость;
    \item простота монтажа;
    \item удобная звездообразная топология;
    \item поддержка питания PoE;
    \item широкая распространённость.
\end{itemize}

Недостатки:
\begin{itemize}
    \item ограничение по длине сегмента;
    \item чувствительность к сильным электромагнитным помехам;
    \item меньшая дальность по сравнению с оптоволокном.
\end{itemize}

\subsubsection{Оптоволокно}

Оптоволокно передаёт данные в виде световых импульсов.

Преимущества:
\begin{itemize}
    \item высокая скорость;
    \item большая дальность;
    \item устойчивость к электромагнитным помехам;
    \item электрическая изоляция;
    \item подходит для магистралей.
\end{itemize}

Недостатки:
\begin{itemize}
    \item более высокая стоимость оборудования;
    \item требовательность к монтажу;
    \item необходимость оптических модулей;
    \item сложность ремонта по сравнению с медным кабелем.
\end{itemize}

\subsubsection{Радиоканал}

Радиоканал используется в Wi-Fi.

Преимущества:
\begin{itemize}
    \item мобильность;
    \item отсутствие необходимости прокладывать кабель к каждому клиенту;
    \item удобство для переносных устройств;
    \item быстрое развёртывание.
\end{itemize}

Недостатки:
\begin{itemize}
    \item помехи;
    \item разделяемая среда;
    \item зависимость от расстояния и препятствий;
    \item необходимость криптографической защиты;
    \item меньшая предсказуемость задержек.
\end{itemize}

\subsection{Итоговые выводы}

Ethernet, Token Ring, FDDI и Wi-Fi представляют разные подходы к построению локальных сетей.

\begin{itemize}
    \item \textbf{Ethernet} эволюционировал от общей коаксиальной шины с CSMA/CD к коммутируемой полнодуплексной архитектуре на витой паре и оптоволокне. Его главные достоинства --- простота, масштабируемость, низкая стоимость и высокая скорость.
    \item \textbf{Token Ring} использует логическое кольцо и маркерный доступ. Его сильная сторона --- детерминированность доступа, но сложность и стоимость привели к вытеснению Ethernet.
    \item \textbf{FDDI} развивает идею маркерного кольца для высокоскоростных и отказоустойчивых оптоволоконных магистралей. Двойное кольцо обеспечивает восстановление при отказах, но технология стала менее актуальной после распространения Fast Ethernet и Gigabit Ethernet.
    \item \textbf{Wi-Fi} обеспечивает беспроводной доступ к локальной сети. Его архитектура строится вокруг точек доступа, BSS и ESS, а доступ к радиосреде организуется с помощью CSMA/CA.
\end{itemize}

Современные локальные сети чаще всего строятся как комбинация:
\begin{itemize}
    \item коммутируемого Ethernet на уровне проводной инфраструктуры;
    \item оптоволоконных магистралей между коммутаторами;
    \item Wi-Fi для мобильного доступа пользователей.
\end{itemize}

\end{document}
